\chapter{系统设计}

\section{总体架构}

本文系统由三部分组成。第一部分是 Python 编排层，负责参数解析、ISA 检测、worker 管理、结果汇总和离线分析入口。第二部分是共享执行核心 \code{injector\_core}，负责调度当前 ISA 后端、写入测试指令、触发执行并将结果输出为统一格式。第三部分是 \code{arch\_riscv.c}、\code{arch\_aarch64.c} 和 \code{arch\_mips.c} 等架构后端，用于封装指令长度、寄存器布局、黑名单规则、跳板与信号语义等 ISA 特定差异。

这样的切分有两个直接收益。其一，前端控制逻辑和后端架构差异被显式隔离，从而避免每支持一个新 ISA 都重写整条流水线。其二，日志协议可以在 C 侧执行器与 Python reader 之间保持稳定，使多 worker 并行和离线聚合都建立在统一帧格式之上。

\section{扫描主循环}

一次扫描的基本流程为：写入待测编码；布置哨兵指令；转移执行流到测试页；捕获异常信号及其上下文；调用 \capstone\ 对编码做反汇编；最后基于执行结果和反汇编结果完成分类，并将结果写入同步文件与最终日志。该主循环本身并不依赖某一特定 ISA，但它要求后端提供一组明确接口，包括测试页准备、跳板进入方式、寄存器沙箱和异常解释。

当前实现同时保留 text 与 raw 两种输出协议。text 输出便于人工阅读和归档，raw 输出面向 \code{sifter.py} 的 reader 线程，能够用固定长度帧传递 worker 标识、编码、反汇编状态、信号号与 \code{si\_code} 等字段。相比完全依赖文本管道，raw 协议更适合多 worker 并发场景，因为它减少了解析歧义，也减轻了 reader 的字符串处理负担。

\section{\memcage\ 路径}

\memcage\ 的核心思想是在用户态构造一个受限执行环境：测试页负责注入待测指令，trap page 和 guard page 吸收大多数寄存器跳转或越界访存，超时处理逻辑负责收回死循环或等待型指令。其直接优势是吞吐更高，因为执行路径主要停留在用户态内部，不必在每次测试时都引入 tracer/tracee 的系统调用切换。

不过，\memcage\ 的高吞吐并不是“免费”的。它至少依赖三类强 ISA 特化设计。第一，哨兵指令必须既能稳定产生可识别后继行为，又不能和常见测试指令发生过强语义干扰。第二，寄存器沙箱需要把大多数通用寄存器预置到安全地址，以降低随机访存、错误返回或寄存器跳转对 harness 的破坏。第三，信号处理逻辑必须结合故障地址或等价上下文，区分“测试指令本身被拒绝”和“测试指令执行成功后命中哨兵”这两类情况。正因如此，本文把 \memcage\ 明确视为高吞吐但高特化的后端。

\placeholderfigure{fig:memcage}{[图 4-1 待补]\\\memcage\ 内存布局\\测试页、trap page、guard page、signal return path}{\memcage\ 路径的内存布局与信号回收流程。终稿将补入真实结构图，用于解释为何该后端吞吐更高，同时也更依赖 ISA 特化 hardening。}

\section{\ptracebackend\ 路径}

\ptracebackend\ 路径通过 tracer/tracee 模式对单条指令执行进行更强的可观察控制。与 \memcage\ 相比，它更容易 bring-up，也更便于调试，因为寄存器读写、程序计数器调整和单步行为都可以通过通用接口显式观察；代价则是 stop/resume 与系统调用引入的额外开销。

从设计选择上看，本文并没有让 \ptracebackend{} 取代 \memcage{}。原因不是 \ptracebackend{} 不够通用，而是两者回答的问题不同。\ptracebackend{} 更适合作为默认 bring-up 路径和调试后端，而 \memcage{} 更适合作为在用户态内追求更高吞吐的强化版本。也正因此，本文的生成器优先产出一个可以接入统一编排框架的基线后端，再把 \memcage\ 相关 hardening 通过 TODO 清单显式暴露出来。

\placeholderfigure{fig:ptrace}{[图 4-2 待补]\\\ptracebackend\ tracer/tracee 单步流程\\tracee 准备、单步、信号回收、状态判定}{\ptracebackend\ 路径的控制流程。该图将用于强调其更强的可观察性与更高的系统调用开销。}

\section{配置驱动生成链路}

为降低新 ISA bring-up 成本，本文设计了配置驱动的后端生成器。其输入是 ISA 描述文件，例如 \code{arch-specs/mips64el.json}；输出则是可编译的基线后端 \code{src/arch\_mips.c} 以及面向人工加固的 \code{MEMCAGE\_TODO} 文档。这个设计有意避免宣称“零人工支持新 ISA”。相反，它明确区分：哪些信息可以安全地由配置推导，哪些部分必须依赖架构知识和后续实测。

在当前实现中，可由配置直接描述的内容包括：\capstone\ 架构与 mode 宏、固定指令长度、基础黑名单、哨兵指令、寄存器组结构，以及 \ptracebackend\ 所需的 PC/GPR 字段名。相比之下，寄存器沙箱策略、异常语义、\code{si\_addr} 解释和特权敏感指令的黑名单收敛，则更适合通过 TODO 清单显式暴露给开发者。这种分工虽然保守，但更符合系统工程中的责任边界。

\section{多 worker 分片与聚合}

除执行路径外，本文的另一项关键设计是多 worker 分片。当前端以 exhaustive 模式并指定 \code{-j N} 时，编码空间会被划分为多个 shard，由独立 worker 分别处理。这样的设计可以在不共享复杂状态的前提下实现并发扩展，同时让每个 worker 的局部状态、最后编码和崩溃计数都能单独记录到 \code{run.json} 中。

与单进程内部多线程方案相比，独立进程更利于故障隔离。即使个别 worker 意外退出，前端也能够根据其分片边界和最后编码恢复局部进度，而不必整体中止全部任务。当前论文并不把 restart 机制包装成单独贡献，但这一工程结构为后续更大规模扫描保留了必要的稳定性空间。

\placeholderfigure{fig:workers}{[图 4-3 可选]\\多 worker 分片、reader、manager 与结果聚合流程}{多 worker 分片与结果聚合流程。该图不是主线必需，但若终稿空间允许，可用于说明前端如何在不理解 ISA 细节的前提下统一聚合多个执行器的输出。}
